\section{Visibility interventions}
\label{sec:visibility}

\subsection{Setting and notation}
\label{sec:setting}

A scan is given by its native voxel centres $v\in\mathbb{R}^3$ in physical coordinates (mm) with intensities $I(v)$ in Hounsfield units (HU), and $p\in\mathbb{R}^3$ denotes a proposal location from the released LUNA16 candidate list~\cite{setio2017luna16}, which we use without modification. For every proposal, the ranker receives a $32^3$ crop centred at $p$ and aligned with the native axes, sampled at $1.2\mm$ by trilinear interpolation, clipped to $[-1000,400]\HU$ and quantized to 8~bits. The networks operate on the central $28^3$ samples, which span $32.4\mm$. All visibility conditions share this crop, so that the input is identical in shape across conditions.

\subsection{Masking before resampling}
\label{sec:mask}

A \emph{visibility condition} is a set $V$ of native voxels that the ranker may observe; all other voxels are replaced by air,
\begin{equation}
 I_V(v)=
 \begin{cases}
 I(v), & v\in V,\\
 -1000\HU, & v\notin V ,
 \end{cases}
 \label{eq:mask}
\end{equation}
and the crop is resampled from $I_V$. Let $\sigma=(\sigma_1,\sigma_2,\sigma_3)$ denote the native voxel spacing of the scan. The value of the crop at a sample point $x\in\mathbb{R}^3$ is the trilinear interpolation
\begin{equation}
 X_V(x)=\sum_{v} w_v(x)\,I_V(v),\quad
 w_v(x)=\prod_{k=1}^{3}\Big[1-\frac{|x_k-v_k|}{\sigma_k}\Big]_+ ,
 \label{eq:trilinear}
\end{equation}
where $[\,\cdot\,]_+=\max(0,\cdot)$. The weights are non-negative and sum to one, and at most eight voxels have $w_v(x)>0$. Masking the native volume rather than the resampled grid is essential: masking after resampling would keep terms $w_v(x)\,I(v)$ of hidden voxels close to the boundary, whereas under Eq.~\eqref{eq:mask} every hidden voxel contributes only $w_v(x)\cdot(-1000\HU)$, a value that does not depend on the hidden intensity. For example, a sample with weight $0.625$ on a visible voxel of $-500\HU$ and $0.375$ on a hidden voxel of $+100\HU$ would read $-275\HU$ if the grid were masked after resampling, but reads $-687.5\HU$, the value obtained with air, under Eq.~\eqref{eq:mask}. For efficiency, the mask is applied to a neighbourhood of each proposal with a margin for interpolation. We verified that this implementation is bit-identical to masking the whole volume, including for rotated, anisotropic and boundary configurations, and that randomizing all hidden values leaves every input unchanged (supplement Sec.~S2). The network receives no explicit visibility channel, so hidden regions appear only as air.

The baseline condition is the native ball
\begin{equation}
B(p)=\{v:\lVert v-p\rVert\le 10\mm\},
\end{equation}
which we compare with the ball together with the band $A(p)$ defined below and with the full field.

\subsection{The exterior band}
\label{sec:band}

Let $d=\lVert \sigma\rVert_2$ denote the diagonal of a native voxel. By Eq.~\eqref{eq:trilinear}, $w_v(q)>0$ requires $|q_k-v_k|<\sigma_k$ on every axis, so any native voxel that contributes to a sample point $q$ satisfies $\lVert v-q\rVert<d$. We place samples at the two radii
\begin{equation}
 \begin{aligned}
 r_N &=10\mm+\max(1.9\mm,\,d+0.05\mm),\\
 r_F &=r_N+1.5\mm ,
 \end{aligned}
 \label{eq:radii}
\end{equation}
For a sample $q$ at distance $r=\lVert q-p\rVert\in\{r_N,r_F\}$ from the proposal, the triangle inequality gives $r-\lVert v-q\rVert\le\lVert v-p\rVert\le r+\lVert v-q\rVert$, and hence every native voxel $v$ contributing to a sample at either radius satisfies
\begin{equation}
 r_N-d\;\leq\;\lVert v-p\rVert\;\leq\;r_F+d .
 \label{eq:support}
\end{equation}
We define the band as exactly this shell of native voxels,
\begin{equation}
A(p)=\{v:\ r_N-d\le\lVert v-p\rVert\le r_F+d\}.
\label{eq:band}
\end{equation}
Since $r_N-d\ge10.05\mm$, the band is disjoint from the ball and separated from it by a thin layer of air. For a spacing of $0.70\times0.70\times1.25\mm$, for example, $d=1.59\mm$, $r_N=11.9\mm$ and $r_F=13.4\mm$, and the band spans $10.3$--$15.0\mm$. Because the radii adapt to the spacing of each scan, Eq.~\eqref{eq:support} holds for any slice thickness.



\subsection{Sparse samples and their native sources}
\label{sec:source_method}

The sparse samples are taken along 96 unit axes $u_1,\dots,u_{96}$ distributed over a hemisphere by a Fibonacci lattice. Sampling both directions of each axis at both radii yields $96\times2\times2=384$ samples $q_{j,\pm,r}=p\pm r\,u_j$ with $r\in\{r_N,r_F\}$. Their values are clipped and normalized and are complemented, for each axis, by a background level computed as the median over all other axes. By construction, $A(p)$ is the envelope of the support of these samples.

The band and the sparse samples differ in two respects: the band exposes many more voxels, and it enters the image network at its input, whereas the samples are read by a separate head. To separate these factors, we introduce a fourth image condition. Let
\begin{equation}
U(p)=\big\{v:\ w_v(q_{j,\pm,r})>0\ \text{for some } j,\ \pm,\ r\big\}
\label{eq:sources}
\end{equation}
be the set of native voxels with strictly positive trilinear weight for at least one of the 384 samples; by Eq.~\eqref{eq:support}, $U(p)\subseteq A(p)$. The \emph{native-source} condition uses $V=B(p)\cup U(p)$ in Eq.~\eqref{eq:mask} and is trained with the same two-scale network and schedule. It thus provides the image network with every native value from which the samples are computed, at the same point of entry as the band, while exposing on average \QuerySources{} exterior voxels per proposal compared with \BandSources{} for the band. Fig.~\ref{fig:visibility} shows the four image conditions, and Algorithm~\ref{alg:band} summarizes the construction of one input for all of them.

\begin{algorithm}[t]
\caption{Construction of one visibility-controlled input.}
\label{alg:band}
\small
\begin{algorithmic}[1]
\Require scan $\{(v,I(v))\}$, voxel spacing $\sigma$, proposal $p$, condition $c\in\{\text{ball},\text{band},\text{sources},\text{full}\}$
\State $d\gets\lVert \sigma\rVert_2$ \Comment{native voxel diagonal}
\State $r_N\gets 10+\max(1.9,\,d+0.05)$;\ \ $r_F\gets r_N+1.5$ \Comment{Eq.~\eqref{eq:radii}}
\State $B\gets\{v:\lVert v-p\rVert\le10\}$
\State $A\gets\{v:r_N-d\le\lVert v-p\rVert\le r_F+d\}$ \Comment{Eq.~\eqref{eq:band}}
\State $U\gets\{v:\ v$ has positive trilinear weight for one of the 384 samples $q_{j,\pm,r}\}$
\State $V\gets B$, $B\cup A$, $B\cup U$ or all voxels, according to $c$
\State $I_V(v)\gets I(v)$ if $v\in V$, else $-1000$ \Comment{Eq.~\eqref{eq:mask}}
\State $X\gets$ trilinear resample of $I_V$ on the $32^3$ grid at $1.2\mm$ around $p$
\State \Return $\mathrm{uint8}(\mathrm{clip}(X,-1000,400))$
\end{algorithmic}
\end{algorithm}

\subsection{Sparse readout heads}
\label{sec:sparse_method}

To assess whether the 384 sample values are sufficient under a suitable readout, we train readout heads on top of a frozen ball-only ranker. Given the logit $s$ of this ranker, a head predicts a correction $h$ from the samples and the clipped logit, and the final score is $s+h$. The last layer of $h$ is initialized to zero, so that each head starts from, and is compared with, the ball-only ranker.

Each axis is represented by a token. A shared encoder maps the two samples on each half-axis, taken relative to the background level, to vectors $a_j,b_j\in\mathbb{R}^c$, where $c$ is the head width, and the token
\begin{equation}
 t_j=\big[\,a_j+b_j,\ a_j\odot b_j,\ |a_j-b_j|\,\big]
 \label{eq:token}
\end{equation}
is invariant to the orientation of the axis, since exchanging $a_j$ and $b_j$ leaves it unchanged. Two attention blocks followed by mean and max pooling aggregate the 96 tokens. In head $m\in\{1,\dots,4\}$ of a block, the attention logit between axes $j$ and $k$ is
\begin{equation}
 \ell^{(m)}_{jk}=\frac{(W^{(m)}_Q t_j)^{\top}(W^{(m)}_K t_k)}{\sqrt{c}}+\beta_m(z_{jk}),
 \label{eq:attention}
\end{equation}
where $W^{(m)}_Q$ and $W^{(m)}_K$ are learned projections and $z_{jk}=|u_j^\top u_k|\in[0,1]$ is the absolute cosine between the two axes. The \emph{direction-aware} variant uses the fixed biases $\beta_1=12(z-1)$ and $\beta_2=35(z-1)$, which favour neighbouring axes at two scales, $\beta_3=-6z$, which favours well-separated axes, and $\beta_4=0$. The \emph{direction-blind} variant sets $\beta_m\equiv0$ in all heads and is otherwise identical, sharing initialization, training batches, objective, checkpoints and parameter count. Tokens and bias matrix are permuted jointly during training, so that direction cannot be inferred from token order. Both variants are trained with widths 48 and 96 (54k and 213k parameters).

We further evaluate a compact head with 873 parameters that pools the axis tokens without attention; besides its initial fit, we refit it on the same frozen scores and data as the attention heads. Four controls are built on this head (Sec.~\ref{sec:compact_result}; supplement Secs.~S2--S3): a single scalar per axis instead of its endpoints, an image embedding in place of the logit, shuffled or removed pairing of opposite half-axes, and the head applied on top of the full-field ranker.

\section{Experimental setup}
\label{sec:evaluation}

\paragraph{Data.} We use the LUNA16 release of LIDC-IDRI~\cite{setio2017luna16,armato2011lidc} with its V2 candidates, reference nodules and excluded findings. Scans are split by patient into a training (356 scans), a selection (44), a development (45) and a calibration (88) set, used for training, checkpoint selection, design choices and threshold calibration, respectively. The evaluation pool consists of the \EvalScans{} scans of subsets 6--9, after removing two scans of a subject that appears twice, and contains \EvalMarks{} proposals and \EvalLesions{} reference lesions, of which at most \ProposedLesions{} have a matching proposal. Before evaluation, the pool was divided by LUNA16 subset into a confirmation half (177 scans) and a replication half (176 scans), which we use to check the stability of the effects. Since earlier studies, including crop studies, used all nodule-bearing scans of this pool, our results are retrospective evaluations on a reused benchmark; all comparisons and the anatomy grouping were therefore fixed before any evaluation score was computed.

\paragraph{Rankers.} Our main ranker is a \emph{two-scale} 3D CNN that processes a central $16^3$ view and a $14^3$ average-pooled view of the $28^3$ input. Each stream consists of a $3^3$ stem with 16 channels followed by three stride-two stages with 32, 64 and 96 channels, and pooled mean and max features are combined by a 96-unit projection (1.28M parameters). To assess transfer, we additionally train a \emph{single-stream} 3D network, which rearranges $2^3$ blocks into eight channels on a $14^3$ grid and processes them with residual stages of 32, 64 and 96 channels (0.68M parameters), and a \emph{three-view} 2D model that reads six central slices per orientation with a shared encoder. All rankers receive the same masked crops.

\paragraph{Implementation details.} All visibility conditions are trained from scratch with an identical schedule: 48 epochs of 160 batches (24 positives and 104 negatives each), AdamW with a learning rate of 0.0018, cosine decay to 12\% and weight decay 0.012, and online hard-negative mining, in which half of the negatives of each batch are drawn from the 16,000 highest-scoring training negatives, refreshed every four epochs. Augmentations (shifts of up to four grid steps, in-plane $90^\circ$ rotations, flips, and intensity gain and offset) are applied after masking. Model selection uses the CPM on the selection set, evaluated every four epochs, with ties resolved in favour of the earlier epoch; the sparse heads follow the same rule over 24 epochs. For the two-scale ranker, we report two training seeds, referred to as \emph{seed~1} and \emph{seed~2}.

\paragraph{Metrics.} Following the release protocol~\cite{setio2017luna16}, a proposal detects a lesion if it lies within the reference radius, only the highest-scoring detection per lesion is counted, and marks on excluded findings are ignored. For a score threshold $\tau$, let $F(\tau)$ be the number of false marks with score at least $\tau$ divided by the number of scans, and let $S(\tau)$ be the fraction of reference lesions whose best matching proposal scores at least $\tau$. Tracing $(F(\tau),S(\tau))$ over $\tau$ gives the free-response ROC (FROC) curve, and $S(r)$ denotes the sensitivity at $r$ false positives (FP) per scan, obtained with the interpolation of the release. The CPM~\cite{niemeijer2011combining,setio2017luna16} averages it over seven rates,
\begin{equation}
\mathrm{CPM}=\tfrac{1}{7}\textstyle\sum_{r\in\{\frac18,\frac14,\frac12,1,2,4,8\}} S(r).
\end{equation}
We additionally report $S_1=S(1)$ and, for each architecture, the fraction of the ball-to-full gap recovered by the band,
\begin{equation}
 \rho=\frac{\mathrm{CPM}_{\text{band}}-\mathrm{CPM}_{\text{ball}}}{\mathrm{CPM}_{\text{full}}-\mathrm{CPM}_{\text{ball}}}.
 \label{eq:recovery}
\end{equation} The release limits each scan to 100 marks, and all comparisons are repeated without this cap (supplement Sec.~S3). Since every model is evaluated at the same FP rate on its own FROC curve, these comparisons measure ranking quality at \emph{matched} operating points; a separate analysis applies thresholds fixed on the calibration set.

\paragraph{Statistical analysis.} All effects are paired differences $\Delta=M(\text{A})-M(\text{B})$ of a metric $M$ (CPM or $S_r$) between two models A and B scored on the same proposals. For $b=1,\dots,2000$, we resample patients with replacement, recompute both FROC curves on the resampled scans and record $\Delta^{(b)}$; the 95\% confidence interval is derived from the distribution of $\{\Delta^{(b)}\}$~\cite{efron1993bootstrap}. These intervals capture uncertainty over patients for fixed trained models, not variation across training runs. The interval of the primary comparison (band versus ball, CPM, seed~1) is adjusted for the three comparisons specified before evaluation, namely band restoration and two contrasts of the compact head (endpoints versus a single scalar, and image embedding versus image score; the first two rows of Fig.~\ref{fig:sparse}b); all other intervals are pointwise.

\paragraph{Anatomy groups.} Reader contours from LIDC-IDRI~\cite{armato2011lidc,fedorov2020lidc} are rasterized on the native grid, and a voxel is labelled as lesion tissue if a strict majority of the linked readers include it. Before any detection scores were computed, each lesion with contours and a matching proposal is said to \emph{extend into the band} if the band of at least one matching proposal contains lesion tissue, so that restoring the band shows part of the lesion itself; otherwise it \emph{does not reach the band}, and the band shows only surrounding tissue. Group sensitivities are computed on the FP axis of the complete proposal stream of each model and normalized by group size, so that comparisons at 1~FP/scan refer to the same global false-positive budget.
